\documentclass[12pt,a4paper]{article}
\usepackage[utf8]{inputenc}
\usepackage[spanish]{babel}
\usepackage{amsmath,amssymb,amsfonts}
\usepackage{geometry}
\usepackage{hyperref}
\geometry{margin=2.5cm}

\title{\textbf{Constantes Matemáticas en Geometría}}
\author{Compendio de Constantes Fundamentales}
\date{\today}

\begin{document}

\maketitle

\section*{Introducción}
La geometría es la rama de las matemáticas que estudia las propiedades, medidas y relaciones de puntos, líneas, ángulos, superficies y sólidos. Las constantes geométricas son fundamentales para determinar áreas, volúmenes, curvaturas y relaciones de escala en el espacio euleriano y no euleriano.

\section{Número Pi ($\pi$)}

\subsection*{1. Significado, Símbolo y Explicación}
Simbolizado con la letra griega $\pi$, representa la razón constante entre la longitud de una circunferencia y su diámetro ($C/d$). Es un número \textbf{trascendental} e \textbf{irracional}, lo que implica que no puede ser raíz de ningún polinomio no nulo con coeficientes enteros. Aparece intrínsecamente en fórmulas de áreas, volúmenes (esferas, cilindros, conos) y ángulos en radianes.

\subsection*{2. Valor Típico en Ingeniería y Ciencia}
En proyectos de ingeniería mecánica, civil y física general:
\begin{equation*}
\pi \approx 3.1416 \quad \text{o} \quad 3.14159265
\end{equation*}

\subsection*{3. Valor Exacto a 15 Decimales}
\begin{equation*}
\pi \approx 3.141592653589793
\end{equation*}

\subsection*{4. Valor Exacto a 100 Decimales}
\begin{align*}
\pi \approx \; & 3.14159265358979323846264338327950288419716939937510582097494459230781... \\
& ...640628620899862803482534211706798215
\end{align*}
\textbf{Margen de error:} Valor matemático puro exacto (sin margen de error).

\vspace{0.8cm}
\hrule
\vspace{0.8cm}

\section{Constante Tau ($\tau$)}

\subsection*{1. Significado, Símbolo y Explicación}
Representada por la letra griega $\tau$, se define como la razón entre la longitud de una circunferencia y su \textbf{radio} ($C/r$). Formalmente es equivalente a $2\pi$. Muchos matemáticos proponen su uso por simplificar la descripción de un giro completo ($1\text{ vuelta} = \tau\text{ radianes}$) y simplificar ecuaciones geométricas y trigonometría de círculos.

\subsection*{2. Valor Típico en Ingeniería y Ciencia}
\begin{equation*}
\tau \approx 6.2832 \quad \text{o} \quad 6.2831853
\end{equation*}

\subsection*{3. Valor Exacto a 15 Decimales}
\begin{equation*}
\tau \approx 6.283185307179586
\end{equation*}

\subsection*{4. Valor Exacto a 100 Decimales}
\begin{align*}
\tau \approx \; & 6.28318530717958647692528676655900576839433879875021164194988918461563... \\
& ...281257241799725606965068423413596430
\end{align*}
\textbf{Margen de error:} Exacto por definición ($\tau = 2\pi$).

\end{document}
